\section{Soundness experiment and bound}
\label{sec:soundness}

\subsection{Why one experiment matters}

A list of individually small probabilities is not yet a security theorem.  The
events must be defined on the same probability space, false acceptance must be
contained in their union, and each event must be charged once.  Otherwise a
bound can omit a branch, count one branch twice, or combine probabilities from
incompatible experiments.

The module \code{V5UnifiedSecurityExperiment} gives the V5 endpoint this
shape.  It uses a finite probability mass function from mathlib rather than a
new probabilistic programming language.  The complete coins may include
random-oracle answers, prover choices, and any other finite randomness needed
by a concrete instantiation.

\begin{experiment}[Work-normalized V5 adversarial experiment]
\label{exp:unified}
For a finite type \(\Omega\), a
\code{WorkNormalizedV5AdversarialExperiment Omega} contains:
\begin{enumerate}[leftmargin=*]
  \item one probability mass function \(\mu\) on \(\Omega\);
  \item a set \(A_{\mathrm{false}}\) of accepted false outcomes;
  \item one record containing all named failure events; and
  \item a pointwise proof that
        \(A_{\mathrm{false}}\) is contained in their unified union.
\end{enumerate}
The law is explicitly work-normalized.  It is not the conditional law of an
attempt after the adversary has completed the grind.
\end{experiment}

Existing accepted-execution reductions can construct this experiment through
\code{ofFinalClassification}.  The more specific constructor
\code{ofReleasedAcceptedExecution} takes the maintained released-execution
classification and its coverage proof.  These constructors prove event
inclusion; they do not create probability bounds for the included events.

\subsection{Exact failure ledger}

An older final ledger has twenty-four entries.  Six belong to the ideal proof
system: the query/final-work miss, four FRI rounds, and relation repair.  The
corrected arithmetic treats their union as one protocol event.  The remaining
eighteen entries describe implementation, primitive, credential, Merkle, and
runtime failures.

Lean proves that the two lists are duplicate-free, disjoint, and concatenate
to the old twenty-four-entry order.  It then defines a nineteen-entry unified
list: one protocol event followed by the eighteen external events.  The
theorem \code{total_unified_failure_eq_total_final_failure} proves equality of
the two union events, not merely an implication.

\begin{table}[p]
\centering
\small
\caption{The nineteen failure kinds in the unified experiment.}
\label{tab:unified-ledger}
\begin{tabularx}{\textwidth}{@{}R{0.07\textwidth}R{0.34\textwidth}Y@{}}
\toprule
No. & Event & Meaning \\
\midrule
1 & Work-normalized protocol event & Union of the query/final-work miss,
four FRI-round events, and relation repair \\
2 & Transcript Rust-to-Lean divergence & Production transcript behavior
differs from the exact Lean driver \\
3 & SHA-256 implementation divergence & Reached SHA-256 call differs from
the specified function \\
4 & SHA-256 collision & Distinct reached inputs have one digest \\
5 & SHA-256 preimage event & The adversary reaches a stated target digest \\
6 & SHA-256 random-oracle event & The ideal-oracle assumption fails in the
used reduction \\
7 & Poseidon2 implementation divergence & Deployed and modeled Poseidon2
differ on a reached input \\
8 & Poseidon2 collision & Distinct reached inputs collide \\
9 & Poseidon2 preimage event & A stated target image is reached \\
10 & Accepted-run relation bridge & Accepted production execution does not
produce the modeled relation success \\
11 & Proof/Merkle opening bridge & Production opening acceptance does not
produce the authenticated forest used by the algebraic checks \\
12 & Victim credential recovery & The adversary learns the victim's secret
credential \\
13 & Rust/state-model mismatch & Production state behavior differs from the
Lean transition \\
14 & System program or address mismatch & Account creation or address
derivation differs from the model \\
15 & Writable-account lock failure & Conflicting writes are not serialized as
assumed \\
16 & Rejection rollback failure & A rejected transaction leaves a visible
partial state change \\
17 & Marker persistence failure & A committed spent marker does not persist \\
18 & Finalized-state observation failure & The observed finalized state does
not match the modeled committed state \\
19 & Close/refund mismatch & Proof close or refund differs from the modeled
recipient and balance \\
\bottomrule
\end{tabularx}
\end{table}

The list is not a claim that all events have nonzero probability.  It is the
complete interface of the present endpoint.  A later source proof can show
that some event is empty or assign it budget zero.  Until then, the event
remains visible.

\subsection{Symbolic union bound}

Let \(C\) be the single protocol event and let
\(E_1,\ldots,E_{18}\) be the external events.  From the pointwise
classification and measure monotonicity, Lean proves
\begin{equation}
  \Prob[A_{\mathrm{false}}]
    \leq \Prob[C] + \sum_{i=1}^{18}\Prob[E_i].
  \label{eq:symbolic-union}
\end{equation}
This is
\code{accepted_false_probability_le_core_plus_external_sum}.  The theorem
uses one PMF and maps the exact ordered list, so there is no separate
independence assumption.  It is an ordinary union bound.

Lean stores the eighteen external bounds in
\code{AssumedExternalSecurityBounds}.  Its fields cover the transcript and
primitives, relation bridge, Merkle bridge, credential recovery, and seven
runtime events.  The theorem \code{external_budget_sum_eq_total} proves that
their sum is the declared external budget.  Therefore
\begin{equation}
  \sum_{i=1}^{18}\Prob[E_i]\leq B_{\mathrm{external}}
\end{equation}
follows from the individual assumptions without introducing a second core
term.

\subsection{Query probability}

The query calculation is one of the fully finite parts of the model.  Each
32-byte squeeze block provides eight little-endian words.  The sampler masks
each word to seventeen bits, scans in order, and keeps only first occurrences.
It accepts after collecting eighteen distinct positions and rejects if the
64-draw supply is exhausted first.

Conditioned on success, Lean proves that every ordered list of eighteen
distinct positions has the same number of preimages.  For a fixed bad set of
size at most \(A\) in a domain of size \(N\), the probability that every
query lies in the set is
\begin{equation}
  \prod_{j=0}^{17}\frac{A-j}{N-j}.
  \label{eq:without-replacement}
\end{equation}
For the released calculation, \(N=131{,}072\) and \(A\leq6{,}082\).
Without conditioning, the event ``sampler succeeds and every output is bad''
is bounded by the same ratio because exhaustion rejects.

Theorems in \code{V5BoundedQuerySamplerUniformity} prove the uniformity of the
bounded first-occurrence sampler.  The fixed-set ratio and its finite
inequality are in \code{V5WithoutReplacementQuerySoundness}.  These results
do not show that SHA-256 outputs are independent uniform words, that the Rust
sampler equals the finite model, or that an invalid FRI word determines the
required bad set.  Those statements occur in other event branches.

\subsection{FRI and relation reductions}

The low-degree part has three logically different steps.

First, the released encoders and distances must meet the hypotheses of the
cited decoding theorem.  Lean checks the field size, point distinctness,
stored evaluation order, degree limits, distances, agreement thresholds,
rounding, and list multiplicities.  The decoding theorem itself is an
external mathematical input.

Second, candidate selection must be causal and coherent.  The released
adaptive extraction theorem follows one initial list member through four
folds or enters a counted round event.  The list is not chosen again after
each fold.  The width-nineteen correlated-agreement argument handles a
challenge-dependent family of committed lanes without adding an unjustified
factor of 240 to the batching event.

Third, the candidate must satisfy the private-spend relation.  The relation
sumcheck theorem bounds the challenge tuples that can repair false scalar
claims.  The candidate/trace reduction lists the remaining ways in which a
false no-witness statement could avoid such a mismatch: four-claim batching,
lane recombination, public-field binding, arithmetic residual extraction, or
hash/Merkle residual extraction.  The production callback and authenticated
opening connections remain separate external events.

The V5 four-way folding protocol is not identical to the published S-two
protocol~\cite{CarmonEtAl2026Stwo}.  The Lean work proves exact local folding,
candidate, and encoder facts for V5.  It does not claim that a theorem for a
different quotient-based or multi-domain protocol applies without the stated
applicability premise.

\subsection{Corrected numerical core}

The corrected arithmetic uses nineteen batching lanes, highest exponent
eighteen, the nonzero extension-field denominator, exactly thirty public-coin
rounds after the initial committed data, and six separately positioned work
checks.  The dominant batching event has about 71 bits after the 37-bit grind
has been completed.  Charging for that grind produces a modeled core of about
100.56 bits.  Lean proves the conservative inequality
\begin{equation}
  B_{\mathrm{core}}
    \leq \frac{7}{10\cdot2^{100}}.
  \label{eq:core-budget}
\end{equation}
The exact theorem is
\code{corrected_selected_release_core_le_seven_tenths}.  Its type retains the
decoding, code-membership, grinding-output, Rust transcript, authenticated
round, Merkle/PCS, Fiat--Shamir, and cited-theorem premises.  Arithmetic does
not discharge those premises.

\subsection{Final theorem}

The external budget is required to satisfy
\begin{equation}
  B_{\mathrm{external}}
    \leq \frac{3}{10\cdot2^{100}}.
  \label{eq:external-budget}
\end{equation}
The structure \code{WorkNormalizedCoreEventConnection} is the remaining
event-level statement that the probability of the exact protocol event under
the experiment law is bounded by the arithmetic quantity
\(B_{\mathrm{core}}\).

\begin{theorem}[Conditional work-normalized soundness]
\label{thm:conditional-soundness}
For a \code{WorkNormalizedV5AdversarialExperiment}, assume:
\begin{enumerate}[leftmargin=*]
  \item its accepted-false event has the proved pointwise classification;
  \item the exact protocol event is connected to the released arithmetic;
  \item the maintained applicability premises for that arithmetic hold;
  \item the protocol budget satisfies \cref{eq:core-budget};
  \item all eighteen external event bounds hold and satisfy
        \cref{eq:external-budget}.
\end{enumerate}
Then
\[
  \Prob[A_{\mathrm{false}}]\leq2^{-100}.
\]
\end{theorem}

Lean proves this as
\code{work_normalized_accepted_false_probability_le_two_pow_neg_100_of_released_core}.
The proof is short because all uncertainty is already represented in the
experiment, event classification, and budget structures.  The theorem is
useful precisely because its remaining assumptions are readable in its type.

\begin{remark}[Raw and work-normalized probability]
Dividing an attack probability by its expected work is useful for comparing
attacks under a query/work budget.  It is not the probability of failure once
the work has already been paid.  Removing the released 37-bit work charge
restores a factor of \(2^{37}\) to the dominant event.  The raw completed-grind
bound is therefore much weaker than \(2^{-100}\).  No result in this paper
describes the deployed verifier as having a raw 100-bit false-acceptance
probability per completed proof.
\end{remark}

\begin{table}[t]
\centering
\small
\caption{Soundness definitions and theorems.}
\label{tab:soundness-anchors}
\begin{tabularx}{\textwidth}{@{}YY@{}}
\toprule
Purpose & Lean item \\
\midrule
One finite law and accepted-false event &
\code{WorkNormalizedV5AdversarialExperiment} \\
Exact old/new ledger equality &
\code{total_unified_failure_eq_total_final_failure} \\
One core plus eighteen external terms &
\code{accepted_false_probability_le_core_plus_external_sum} \\
Exact external budget sum &
\code{external_budget_sum_eq_total} \\
Bounded distinct-query uniformity &
\code{AspisV5BoundedQuerySamplerUniformity.deployed_q18_conditioned_uniform} \\
Released four-round extraction &
\code{AspisV5FriReleasedAdaptiveExtraction.accepted_ideal_fri_extracts_initial_candidate_or_counted} \\
70-percent core budget &
\code{AspisV5HundredBitSecurityMargin.corrected_selected_release_core_le_seven_tenths} \\
Final 100-bit implication &
\code{work_normalized_accepted_false_probability_le_two_pow_neg_100_of_released_core} \\
\bottomrule
\end{tabularx}
\end{table}
